% Chapter 2 -- Rubric: PO1 (1.1), PO2 (research literature), feeds PO3 and PO11
\chapter{Preliminary Research}
\label{ch:research}

\section{Literature Review}
\label{sec:literature}

Computational analysis of persuasive, multi-turn conversational interaction sits at the convergence of behavioral psychology, natural language processing, and statistical machine learning. This section synthesizes the foundational and recent literature across these areas, highlighting how prior works inform this investigation and identifying the theoretical and empirical gaps this capstone addresses.

\subsection{Persuasion Theory and Rhetorical Frameworks}
The systematic study of interpersonal persuasion originates in classical rhetoric and social psychology. Cialdini's seminal framework of influence \cite{cialdini2001influence} categorizes compliance-gaining mechanisms into core principles: reciprocity, commitment and consistency, social proof, authority, liking, and scarcity. In sequential negotiation, classic experimental psychology established specific multi-stage tactics. Freedman and Fraser \cite{freedman1966compliance} demonstrated the \textit{foot-in-the-door} phenomenon, whereby inducing compliance with an initial trivial request substantially increases adherence to a subsequent larger request. Conversely, Cialdini et al. \cite{cialdini1975reciprocal} formulated the \textit{door-in-the-face} technique, wherein an extreme, predictably rejected opening ask creates psychological pressure for the decider to accept a smaller concessionary follow-up. While these principles are universally acknowledged in social science, their formalization into high-coverage computational taxonomies for natural language dialogues has remained limited, often collapsing nuanced rhetorical moves into broad, coarse-grained dialog act labels.

\subsection{Persuasion Corpora and Conversational Strategy Annotation}
The release of the \textit{Persuasion for Good} (P4G) corpus by Wang et al. \cite{wang2019persuasion} marked a major milestone in empirical persuasion modeling. P4G contains 1,017 multi-turn online conversations where an assigned Persuader attempts to convince a Persuadee to donate part of their \$2.00 task earnings to the charity \textit{Save the Children}. Wang et al. annotated a subset of 300 dialogues with a taxonomy of 10 persuasion strategies. Subsequent works expanded on this foundation: Saha et al. \cite{saha2021towards} incorporated sentiment and emotion dynamics into dialog act recognition, while Tian et al. \cite{tian2020towards} explored strategy-aware response generation. Chen and Yang \cite{chen2021persuasion} published a comprehensive survey synthesizing conversational persuasion systems, emphasizing the need for hierarchical representations. More recently, Chen and Xiao \cite{chen2025tactics} analyzed persuasion tactics in commercial multi-party chats, while Petrova et al. \cite{petrova2026revisiting} explored hierarchical re-annotation of P4G. Despite these advances, existing benchmarks frequently suffer from small annotated sample sizes, flat non-hierarchical label spaces, and a failure to address severe label sparsity in the long tail.

\subsection{Transformer Encoders and Dialogue Language Models}
The evolution of contextual language representation—from BERT \cite{devlin2019bert} to RoBERTa \cite{liu2019roberta} and DeBERTa \cite{he2020deberta}—revolutionized NLP by pre-training deep bidirectional representations on vast corpora. RoBERTa optimized BERT's pre-training regime by removing the next-sentence prediction objective and training on longer sequences with dynamic masking. DeBERTa and DeBERTa-v3 \cite{he2021debertav3} advanced this paradigm through disentangled attention (representing tokens by separate content and relative position vectors) and enhanced mask decoder pre-training with ELECTRA-style \cite{clark2020electra} discriminator losses. In dialogue modeling, Wu et al. \cite{wu2020todbert} introduced TOD-BERT, pre-training on diverse task-oriented dialogue datasets to model turn-level conversational semantics. However, when applied naively to multi-turn asymmetric dialogue, these monolithic sequence encoders face rigid context length limits (256 or 512 tokens), silently dropping terminal turns under default right-truncation.

\subsection{Hierarchical and Role-Aware Dialogue Networks}
To process extended multi-turn interactions without truncation, hierarchical neural architectures decouple token-level encoding from turn-level discourse modeling. Li et al. \cite{li2020hitrans} introduced HiTrans, employing a dual-transformer hierarchy where a lower transformer encodes intra-utterance words and an upper transformer models inter-utterance discourse relations. In conversational emotion recognition, Majumder et al. \cite{majumder2019dialoguernn} introduced DialogueRNN, tracking individual speaker party states using recurrent neural networks (GRU \cite{cho2014learning}) to preserve speaker identity and interpersonal context. While DialogueRNN demonstrated the power of per-party tracking for emotional dynamics, asymmetric outcome prediction requires modeling not merely independent speaker states, but the explicit \textit{directional influence} of the arguer's moves upon the decider's internal commitment.

\subsection{Imbalanced Learning and Loss Formulations}
Multi-label persuasion strategy attribution and binary commitment classification both exhibit extreme class imbalance. Conventional cross-entropy loss causes gradients to be overwhelmed by easy negative examples or majority classes. Lin et al. \cite{lin2017focal} introduced Focal Loss, applying a modulating factor $(1 - p_t)^\gamma$ to down-weight easy examples and focus gradient updates on hard instances. Cui et al. \cite{cui2019class} developed Class-Balanced Focal Loss, weighting the loss inversely by the effective number of samples $(1 - \beta) / (1 - \beta^{n_y})$ rather than raw sample frequency, preventing gradient explosion on extremely rare classes. For multi-label settings where 90\%+ of label matrices are zero, Ben-Baruch et al. \cite{benbaruch2021asymmetric} formulated Asymmetric Loss (ASL), dynamically decoupling the focusing parameters $\gamma_+$ and $\gamma_-$ and performing hard probability threshold shifting on negative samples. Furthermore, Khosla et al. \cite{khosla2020supervised} proposed Supervised Contrastive Loss (SupCon), pulling normalized representations of the same class together in latent space while pushing different classes apart, substantially improving representation clustering for minority categories.

\subsection{Ensembling and Stacking Generalization}
Ensemble learning leverages model diversity to minimize generalization variance. Wolpert \cite{wolpert1992stacked} established the mathematical foundation of Stacked Generalization, training meta-classifiers on the cross-validated out-of-fold predictions of heterogeneous base models. In practical applications, gradient-boosted decision trees such as XGBoost \cite{chen2016xgboost} are widely employed as stacking decision layers due to their ability to capture non-linear feature interactions and resist overfitting. Additionally, probability calibration techniques, originating in Platt scaling \cite{platt1999probabilistic}, map uncalibrated logits into true empirical posterior probabilities, enabling smooth mathematical blending (e.g., power-mean ensembling) that often outperforms greedy discrete selection.

\subsection{Negotiation Benchmarks and Shortcut Learning}
Early computational negotiation frameworks, such as the \textit{Deal or No Deal} benchmark by Lewis et al. \cite{lewis2017deal} and the modular negotiation models of He et al. \cite{he2018decoupling}, examined multi-issue bargaining over item allocations. However, recent critical analyses by Geirhos et al. \cite{geirhos2020shortcut} and Gururangan et al. \cite{gururangan2018annotation} demonstrate that neural models trained on complex textual benchmarks frequently exploit superficial \textit{shortcut learning} and dataset artifacts. In conversational intent prediction, models often achieve artificially high accuracy simply by detecting formulaic closing tokens (e.g., `\texttt{sounds good'', }`deal'') in the final turn, completely bypassing the conversational trajectory. Establishing whether an architecture genuinely models trajectory dynamics requires rigorous adversarial evaluation and input-ablation controls.

\subsection{Synthesis of Research Gaps}
\Cref{tab:lit-gap} synthesizes the key open gaps in the existing literature that this project systematically resolves.

\begin{table}[htbp]
  \centering
  \small
  \caption{Synthesis of research literature and identified gaps}
  \label{tab:lit-gap}
  \begin{tabularx}{\textwidth}{l>{\hsize=0.9\hsize}Y>{\hsize=1.1\hsize}Y}
    \toprule
    \textbf{Research Area} & \textbf{State of the Art / Established Findings} & \textbf{Open Gap Addressed in this Project} \\
    \midrule
    \textbf{Persuasion Modeling} & Coarse-grained 10-class taxonomies on small subsets (300 dialogues) \cite{wang2019persuasion}. & Full 10,600-turn multi-label dataset with an 11-category, 41-strategy hierarchical taxonomy. \\
    \midrule
    \textbf{Intention Grounding} & Trusting recorded numerical corpus metadata (\$0 vs \$10). & Proving textual divergence; conducting 3-pass human re-annotation to create strict commitment labels. \\
    \midrule
    \textbf{Dialogue Architecture} & Flat transformers (256 tokens) or mixed-sequence hierarchical transformers (HiTrans \cite{li2020hitrans}). & Originating \textit{Antiphony}: role-pure dual streams with directional cross-attention ($Q_{\mathcal{D}}, K_{\mathcal{P}}, V_{\mathcal{P}}$) and trajectory GRU. \\
    \midrule
    \textbf{Multi-Label Extreme Skew} & Standard multi-label binary cross-entropy suffering long-tail collapse. & Combining Asymmetric Loss, Class-Balanced Focal Loss, and calibrated power-mean ensemble blending across 19 models. \\
    \midrule
    \textbf{Negotiation Evaluation} & Clean synthetic allocations \cite{lewis2017deal} vulnerable to closing-turn shortcut learning. & Generating 700-dialogue adversarial reversal benchmark; proving Antiphony widens sufficiency gap ($\Delta_{\text{suff}} = +0.0384$). \\
    \bottomrule
  \end{tabularx}
\end{table}

\section{Existing Solutions and Technologies}
\label{sec:related}

To contextualize the technical positioning of this work, we survey the primary categories of existing technologies currently deployed in commercial dialogue analysis and conversational commerce:

\begin{enumerate}
  \item \textbf{Monolithic Transformer Fine-Tuning:} Industrial practice frequently applies standard pre-trained language models (RoBERTa, DeBERTa) directly to flattened dialogue text. While computationally straightforward to implement using libraries such as Hugging Face Transformers \cite{wolf2020transformers}, these models truncate long dialogues at 256 or 512 tokens. As demonstrated empirically in Chapter 5, default right-truncation cuts off the closing exchanges in $30.4\%$ of dialogues, completely blinding the model to the final decision.
  \item \textbf{Prompted General-Purpose Large Language Models (LLMs):} Commercial solutions increasingly query proprietary frontier LLMs (e.g., GPT-4, Claude-3.5) via zero-shot or few-shot prompts to extract customer intent and conversational sentiment. While highly flexible, prompted LLMs present critical operational limitations: prohibitive inference latency ($>1{,}000$\,ms per dialogue), massive recurring financial costs, privacy violations resulting from sending proprietary customer transcripts to third-party endpoints, and severe non-determinism.
  \item \textbf{Task-Oriented Slot-Filling Dialogue Systems:} Traditional conversational frameworks (e.g., Rasa, Google Dialogflow) model dialogue state through predefined intent hierarchies and entity slot-filling. While effective for rigid transactional booking (e.g., flight reservations), they completely lack the capacity to model multi-label rhetorical strategies, social rapport, emotional appeals, and fluid non-monotonic negotiation dynamics.
  \item \textbf{Classical Bag-of-Words and TF-IDF Baselines:} Lightweight production systems often rely on TF-IDF feature extraction combined with logistic regression or gradient-boosted decision trees. While fast and interpretable, our experiments reveal that bag-of-words models degrade significantly when exposed to extended conversational context (dropping from $0.575$ to $0.419$ micro-F1 across 5 turns), lacking deep semantic and contextual reasoning.
\end{enumerate}

\textbf{The Identified Technology Gap:} The industry lacks an efficient, local, privacy-preserving discriminative architecture that captures asymmetric conversational dynamics over long multi-turn horizons without arbitrary context truncation, while remaining trainable on standard commodity GPU hardware.

\section{Market and User Research}
\label{sec:market}

\subsection{The Rise of Conversational Social Commerce}
In emerging markets, particularly across South Asia and Southeast Asia, conversational commerce has outpaced traditional web storefronts. In Bangladesh, thousands of micro-, small, and medium enterprises (MSMEs) operate almost entirely across social media chat channels: Facebook Messenger, WhatsApp Business, and Instagram Direct. Transactions do not occur through standardized shopping carts; instead, the entire customer journey—product inquiry, authenticity verification, price negotiation, delivery coordination, and payment commitment—takes place inside unscripted, natural language chat transcripts.

\subsection{User Personas and Behavioral Dynamics}
Qualitative analysis of real marketplace chat patterns reveals key consumer and merchant behavioral archetypes that drive transactional outcomes:
\begin{itemize}
  \item \textbf{Consumer Friction and Skepticism:} Due to prevalent online shopping fraud, consumers display acute risk aversion. Buyers frequently demand physical shop addresses, unboxing photographs, and real-time serial number verification. Prepayment via mobile financial services (MFS such as bKash and Nagad) is frequently contested; buyers overwhelmingly prefer Cash-on-Delivery (COD) with inspection privileges.
  \item \textbf{Bargaining Culture:} Price haggling is deeply culturally ingrained. Buyers expect merchants to offer discretionary price cuts, promotional freebies, or delivery fee waivers. Conversations frequently stall when merchants maintain rigid pricing policies.
  \item \textbf{Merchant Administrative Burden:} Independent page admins and shopkeepers manage hundreds of active chat inquiries simultaneously. Admins report high operational fatigue caused by ``ghosting'' (buyers conducting exhaustive inquiries and abruptly disengaging) and struggle to prioritize leads that have genuine, high-probability buying intent.
\end{itemize}

\subsection{Industry Consultations and Field Observations}
Informal discussions with local social commerce merchants and digital marketing practitioners \pending{Merchant and user interview records} confirm an urgent demand for automated analytical tools. Specifically, merchants express strong interest in automated systems capable of: (i) triaging high-intent leads in real-time, (ii) auditing sales representative dialogue quality against proven persuasive strategies, and (iii) detecting transactional friction before customer disengagement occurs.

\section{Feasibility Study}
\label{sec:feasibility}

We conducted an exhaustive feasibility assessment across four primary dimensions:
\begin{enumerate}
  \item \textbf{Technical Feasibility:} The proposed approach builds upon proven foundational technologies: PyTorch deep learning frameworks, pre-trained Hugging Face transformer checkpoints, and established optimization libraries. Technical feasibility is substantiated by the comprehensive experimental benchmarks presented in Chapter 5, demonstrating that the developed Antiphony architecture achieves state-of-the-art Macro-F1 ($>0.86$) and trains reliably within 3 to 4 GPU hours.
  \item \textbf{Economic Feasibility:} The entire empirical development lifecycle was deliberately designed to operate within zero-cost compute constraints. All models were trained, evaluated, and ablated on free-tier cloud GPU accelerators (NVIDIA Tesla T4 with 16\,GB VRAM on Kaggle). The synthetic benchmark generation pipeline leveraged free-tier and low-cost API credits \pending{Exact API spending}, demonstrating that advanced conversational NLP research can be successfully executed without institutional cluster budgets.
  \item \textbf{Operational Feasibility:} The trained discriminative classifiers are compact ($<500$\,MB model checkpoints for base encoders) and execute inference on standard server GPUs without specialized distributed serving hardware. They integrate cleanly into offline batch-processing analytics pipelines for post-call audits and customer sentiment tracking.
  \item \textbf{Legal and Regulatory Feasibility:} The research strictly respects intellectual property and privacy mandates. The P4G corpus is an established academic benchmark with open research permissions. No proprietary or private user chat logs were harvested; cross-domain commercial transfer was evaluated exclusively on a high-fidelity synthetic benchmark generated under strict ethical guidelines.
\end{enumerate}
In conclusion, the project demonstrates total feasibility across all four operational dimensions.
